\section{Experiment design}

\subsection{Participants}

 We conducted a quantitative and qualitative study using 3 external participants.
 These participants are P0 (Age 21), P1 (Age 26) and P2 (Age 25) respectively.
 Consent to the usage of data was acquired by having the participants fill in a Google Form.
 In addition, the participants were informed that they could withdraw from the study at any point.

\subsection{Environment}

The participants were seated in a small study cubicle with a table and a chair.
Any distraction and noise was minimized.

\subsection{Equipment}

We have performed a visual-latency study using a laptop screen and webcam.
The setup gives a live-feed of the participants hands and phone behind the laptop screen,
out of direct sight. The exact details of our setup are as follows:

\begin{itemize}
	\item A medium-grade laptop with a 144 Hz refresh rate and a Nvidia GeForce RTX 3050 Laptop GPU.
	\item A BlueBuilt Ultra HD 4K Webcam that supports 60 FPS with a 1080p resolution.
		We favored a higher resolution over a higher quality to minimize any transmission latency that could
		interfere with the experiment.
	\item OBS Studio to create the software setup of reading the webcam feed,
		positioning and scaling it on the screen and applying a visual delay in milliseconds.
	\item The users own phone, with any form of auto-correction and spell-assist disabled.
		This ensured that the users were most comfortable with their typing setup.
\end{itemize}

\subsection{Tasks to perform, instructions given, ...}

App: MonkeyType

Fixed amount of words, 50 words per run, randomly generated from a set of 1000 words (english 1k MonkeyType test with 50-word text).

\subsection{Within/between-subjects design}

Within subject design

\subsection{Order (for within-subjects design) and other mappings (for both)}

Each participant is given two runs of the excersize with zero latency as a practice.

Then they are given 5 runs of the excersize with different delays 
(0ms, 25ms, 50ms, 75ms, 100ms) in a random order.

\subsection{Other considerations, if any}

Participants complete the test on their own smartphone,
with their own typing style and keyboard layout. This
is to eliminate the confounding factor of unfamiliarity
with a new keyboard layout.
